\section*{Bab \ref{ch:b2:magmatic-series} --- Seri Magmatik dan Tatanan Tektonik}

\begin{solution}{exo:b2:magmatic-series:1}
(a) $\ce{SiO2} = 49.60$, alkali $2.22 + 0.51 = 2.73$: di antara 45 dan 52 dan di bawah 5,
\emph{basal} (B). (b) 58.5 dan $3.9 + 1.6 = 5.5$: di antara 57 dan 63; batas
atas di sana berjalan dari (57, 5.9) ke (63, 7.0), yaitu $5.9 + 1.1\times1.5/6 = 6.18$
untuk 58.5; 5.5 berada di bawahnya: \emph{andesit} (O2). (c) 72.0 dan $4.0 + 4.2 = 8.2$:
di seberang garis vertikal pada 69 di atas 8 dan di atas garis dari (69, 8) ke (77, 0), yang
bernilai 5.0 pada 72: \emph{riolit} (R). (d) 47.0 dan $4.5 + 2.5 = 7.0$: pada 47, medan U1
terletak di antara garis dari (45, 5) ke (49.4, 7.3), yaitu $5 + 2.3\times2/4.4 = 6.05$, dan
garis dari (45, 9.4) ke (49.4, 7.3), yaitu $9.4 - 2.1\times2/4.4 = 8.45$. Batuan itu
berada di U1: \emph{tefrit} jika normanya mengandung olivin kurang dari \qty{10}{\percent},
\emph{basanit} jika lebih.
\end{solution}

\begin{solution}{exo:b2:magmatic-series:2}
(a) Jumlah oksidanya 98.40 dan jumlah bersama volatil 100.00;
\ce{H2O} $1.40 < 2$ dan \ce{CO2} $0.20 < 0.5$: analisis itu dapat diterima. (b) Faktor
$100/98.40 = 1.01626$: \ce{SiO2} 52.44, \ce{TiO2} 1.22, \ce{Al2O3} 16.36,
\ce{FeO_T} 9.96, \ce{MgO} 6.71, \ce{CaO} 8.33, \ce{Na2O} 3.56, \ce{K2O} 1.42
(jumlah 100.00); alkali $3.56 + 1.42 = 4.98$. (c) Sebelum: 51.60 dan 4.90, sedikit di bawah
batas basal--andesit basaltik pada 52: \emph{basal}. Sesudah: 52.44 dan 4.98; pada
52.44, batas atas O1, dari (52, 5) ke (57, 5.9), bernilai $5 + 0.9\times0.44/5
= 5.08$: \emph{andesit basaltik} (O1). Perhitungan ulang memindahkan batuan itu melintasi sebuah
batas: nama di dekat suatu garis bergantung padanya. (d) Pada 52.44, garis K-rendah bernilai $0.3 +
0.045\times4.44 = 0.50$ dan garis K-tinggi $1.2 + 0.095\times4.44 = 1.62$; 1.42
terletak di antaranya: \emph{K-sedang}.
\end{solution}

\begin{solution}{exo:b2:magmatic-series:3}
(a) $Q + A + P = 90$: $Q = 25/90 = \qty{27.8}{\percent}$ dan $100\,P/(A + P) =
100\times35/65 = 53.8$: medan 3b, \emph{monzogranit}, sebuah granit. Biotit dan
hornblende tidak masuk ke dalam jumlah QAPF. (b) $Q + A + P = 65$: $Q = 5/65 =
\qty{7.7}{\percent}$ (di antara 5 dan 20) dan $100\times55/60 = 91.7$ (di atas 90): medan
10*, diorit kuarsa, gabro kuarsa, atau anortosit kuarsa. \omterm{def:b2:magmatic-series:qapf}{Diagram QAPF} saja tidak
memilih di antaranya; plagioklas dan mineral mafiknya (di sini
\qty{35}{\percent} klinopiroksen) yang menentukan.
\end{solution}

\begin{solution}{exo:b2:magmatic-series:4}
Garis K-rendah adalah $0.3 + 0.045(\ce{SiO2} - 48)$, garis K-tinggi $1.2 +
0.095(\ce{SiO2} - 48)$. (a) Pada 50: 0.39 dan 1.39; $0.40$ sedikit di atas garis bawah:
\emph{K-sedang}. (b) Pada 56: 0.66 dan 1.96: \emph{K-sedang}. (c) Pada 62: 0.93 dan 2.53;
$2.90 > 2.53$: \emph{K-tinggi}. (d) Pada 66: 1.11 dan 2.91; $1.00 < 1.11$: \emph{K-rendah}.
\end{solution}

\begin{solution}{exo:b2:magmatic-series:5}
(a) Basal: $A + F + M = 3.2 + 14.5 + 4.5 = 22.2$; $A = 14.4$, $F = 65.3$, $M = 20.3$.
Andesit: $5.05 + 6.1 + 3.8 = 14.95$; $A = 33.8$, $F = 40.8$, $M = 25.4$. (b) Garis itu
tidak pernah melampaui $F = 54.5$; basal, dengan $F = 65.3$, terletak di sisi F
berapa pun $A$-nya: \emph{toleiitik}, ferobasal dari tren Fenner. Untuk
andesit, pada $A = 33.8$ garis itu terletak di antara (49.6, 35.0) dan (52.5, 30.0), pada $F =
49.6 + 2.9\times1.2/5 = 50.3$; $40.8 < 50.3$: \emph{kalk-alkali}. (c) Magnetit,
yang sejak awal membawa besi feri keluar dari \omterm{def:b2:magma-emplacement:magma}{magma} yang teroksidasi dan mencegah $F/M$
naik secepat pada toleiit.
\end{solution}

\begin{solution}{exo:b2:magmatic-series:6}
MORB: $\mathrm{Nb_N} = 4/0.658 = 6.08$, $\mathrm{La_N} = 4.12/0.648 = 6.36$, rasio
0.956. \omterm{def:b2:elements-in-the-earth:morb-oib}{Basal pulau samudra}: $38/0.658 = 57.8$, $30.8/0.648 = 47.5$, rasio 1.215. Basal
busur: $3/0.658 = 4.56$, $6.1/0.648 = 9.41$, rasio 0.484. Hanya basal busur yang memiliki
anomali niobium, yaitu lembah pada Nb di diagram laba-laba: niobium terkuras dibandingkan dengan
tetangganya, La: \omterm{def:b2:magma-emplacement:magma}{magma} busur menerima unsur-unsur dari slab, tetapi hanya sedikit
niobium. MORB hampir datar; \omterm{def:b2:elements-in-the-earth:morb-oib}{basal pulau samudra} lebih kaya
akan Nb daripada La.
\end{solution}

\begin{solution}{exo:b2:magmatic-series:7}
(a) $\mathrm{Ti} = 1.6\times0.5993\times10^4 = \qty{9589}{ppm}$; Ti/100 = 95.9, Zr
110, $3\mathrm Y = 96$; jumlah 301.9; koordinat 31.8, 36.4, 31.8: medan B. (b) Ti
16\,780, jadi 167.8, 220, 87; jumlah 474.8; 35.3, 46.3, 18.3: medan D. (c) Ti 4794, jadi
47.9, 50, 60; jumlah 157.9; 30.3, 31.7, 38.0: medan B. Basal (b) adalah basal dalam-lempeng
(dari pulau samudra atau benua); (a) dan (c), di medan B, mungkin MORB, toleiit
busur, atau basal kalk-alkali: diagram itu tidak dapat memutuskan, dan diperlukan diagram lain (diagram
laba-laba, anomali Nb).
\end{solution}

\begin{solution}{exo:b2:magmatic-series:8}
(a) MORB $142/164 = 0.866$, \omterm{def:b2:elements-in-the-earth:morb-oib}{basal pulau samudra} $4/173 = 0.023$, basal busur $101/201 =
0.502$. (b) Dengan peluang awal yang sama, aturan Bayes memberikan $0.866/(0.866 + 0.023 + 0.502) =
0.866/1.391 = 0.623$. (c) Bahwa basal itu sangat kecil kemungkinannya merupakan \omterm{def:b2:elements-in-the-earth:morb-oib}{basal pulau samudra}
(peluang $0.023/1.391 = 0.017$), tetapi tidak dapat menentukan apakah basal itu MORB atau basal busur:
peluangnya kira-kira tiga banding dua.
\end{solution}

\begin{solution}{exo:b2:magmatic-series:9}
Mol per \qty{100}{g}: \ce{Na2O} $5.5/61.98 = 0.08874$, \ce{K2O} $4.5/94.20 =
0.04777$, \ce{Al2O3} $10.5/101.96 = 0.10298$. (a) $(0.08874 + 0.04777)/0.10298 =
0.13651/0.10298 = 1.326 > 1$: riolit itu \emph{peralkali}. (b) Feldspar dapat
menampung $0.10298$ dari $0.13651$ mol oksida alkali; kelebihannya, $0.03353$, adalah
$0.03353/0.13651 = 0.246$, seperempat alkali, yang harus masuk ke mineral lain.
\end{solution}

\begin{solution}{exo:b2:magmatic-series:10}
(a) $s = (\lambda_F - \lambda_M)/(\lambda_F - \lambda_A)$: untuk $\lambda_F = 0.5$, $(0.5
- 0.8)/(0.5 - 0.4) = -3.0$; untuk $0.7$, $(-0.1)/(0.3) = -0.333$. (b) Penyelesaian
$s(\lambda_F - 0.4) = \lambda_F - 0.8$ memberikan $\lambda_F = (0.4s - 0.8)/(s - 1)$:
Islandia, $s = -1.10$: $(-0.44 - 0.8)/(-2.10) = 0.590$; Cascade, $s = -0.72$:
$(-0.288 - 0.8)/(-1.72) = 0.633$. (c) \omterm{def:b2:magma-emplacement:magma}{Magma} Cascade, dengan
$\lambda_F$ yang lebih besar, kehilangan besinya lebih cepat: magnetit mengeluarkannya. (d) $\ln(F/M)$ bertambah dengan
laju $\lambda_M - \lambda_F$ dan $\ln(A/F)$ dengan laju $\lambda_F - \lambda_A$; jadi untuk
setiap satuan yang diperoleh $\ln(A/F)$, $\ln(F/M)$ bertambah $-s$: 1.10 di Islandia, 0.72 di
Cascade. $F/M$ naik di sepanjang kedua tren, tetapi lebih cepat pada tren toleiitik, yaitu tren
Fenner.
\end{solution}

\begin{solution}{exo:b2:magmatic-series:11}
(a) BHVO-2, mol per \qty{100}{g}: \ce{SiO2} 0.82557, \ce{Al2O3} 0.13182, \ce{FeO}
0.15521, \ce{MgO} 0.18007, \ce{CaO} 0.20328, \ce{Na2O} 0.03580, \ce{K2O} 0.00545.
Ortoklas 0.00545 dan albit 0.03580; sisa \ce{Al2O3} $0.13182 - 0.04125 = 0.09057$,
seluruhnya ke dalam anortit bersama \ce{CaO} dalam jumlah yang sama; sisa \ce{CaO} $0.11271$ ke dalam diopsid;
sisa $(\ce{Mg,Fe})\ce{O}$ $0.33528 - 0.11271 = 0.22257$ ke dalam hipersten. Silika
yang diperlukan: $6\times0.00545 + 6\times0.03580 + 2\times0.09057 + 2\times0.11271 + 0.22257 =
0.87663$; yang tersedia 0.82557; kurang 0.05106. Setiap $(\ce{Mg,Fe})\ce{O}$ yang diubah dari
hipersten menjadi olivin menghemat setengah silika, paling banyak $0.22257/2 = 0.11129$: mengubah
$0.10212$ darinya sudah cukup, dan $0.12045$ tetap dalam hipersten. Normanya mengandung olivin
dan hipersten, tanpa kuarsa maupun nefelin: batuan itu jenuh silika dalam
arti definisi tadi, sebuah \emph{toleiit olivin} bernorma hipersten, dari
\omterm{def:b2:magmatic-series:series}{seri toleiitik}, seperti yang diletuskan pada tahap perisai gunung api Hawaii. (b) Mol: 0.75732, 0.14712, 0.16008, 0.21092, 0.17832, 0.06131,
0.01592. Ortoklas 0.01592, albit 0.06131, anortit 0.06988, diopsid 0.10843, hipersten $0.37100 - 0.10843 =
0.26257$; silika yang diperlukan $0.09552 + 0.36786 + 0.13976 + 0.21686 + 0.26257 = 1.08257$;
yang tersedia 0.75732; kurang 0.32525. Mengubah seluruh hipersten menjadi olivin menghemat
0.13129, sehingga kekurangannya tinggal 0.19396; setiap mol \ce{Na2O} yang diubah dari albit menjadi nefelin
membebaskan 4 \ce{SiO2}, jadi $0.19396/4 = 0.04849$
dari 0.06131 mol \ce{Na2O} albit menjadi nefelin. Ada nefelin dalam normanya: batuan itu
kurang jenuh silika, sebuah \emph{basal alkali}, dari \omterm{def:b2:magmatic-series:series}{seri alkali}.
\end{solution}

\begin{solution}{exo:b2:magmatic-series:12}
(a) Basal busur $116/141 = 0.823$, MORB $21/101 = 0.208$, \omterm{def:b2:elements-in-the-earth:morb-oib}{basal pulau samudra}
$3/137 = 0.022$. (b) $0.823/(0.823 + 0.208 + 0.022) = 0.823/1.053 = 0.782$.
(c) Pembilang $0.1\times0.823 = 0.0823$; penyebut $0.0823 + 0.6\times0.208 +
0.3\times0.022 = 0.0823 + 0.1248 + 0.0066 = 0.2137$; peluangnya $0.385$. (d) Uji
yang sama pada batuan yang sama memberikan 0.78 atau 0.39 bergantung pada seberapa lazim setiap tatanan
di tempat batuan itu ditemukan: \omterm{def:b2:magmatic-series:discrimination}{diagram diskriminasi} memberikan kemungkinan, dan
konteks geologi harus melengkapi sisanya.
\end{solution}

\begin{solution}{pb:b2:magmatic-series:1}
\noindent\textbf{Bagian I.}

\textbf{1.} Oksida: lava 1 99.17, lava 2 98.65, lava 3 98.20; bersama volatil
99.97, 99.80, dan 99.45, semuanya mendekati 100. Air (0.70, 1.10, 1.20) di bawah
\qty{2}{\percent} dan karbon dioksida (0.10, 0.05, 0.05) di bawah \qty{0.5}{\percent}: ketiganya
dapat diterima.

\textbf{2.} Faktor $100/99.17 = 1.00837$: \ce{SiO2} 50.82, \ce{Na2O} 2.92, \ce{K2O}
0.63, alkali 3.55.

\textbf{3.} Lava 2: $100/98.65 = 1.01368$, \ce{SiO2} 59.40, \ce{Na2O} 3.95, \ce{K2O}
1.17, alkali 5.12. Lava 3: $100/98.20 = 1.01833$, \ce{SiO2} 66.50, \ce{Na2O} 4.48,
\ce{K2O} 1.88, alkali 6.36.

\textbf{4.} Lava 1 (50.82, 3.55): basal. Lava 2 (59.40, 5.12; batas atas O2
di sana bernilai $5.9 + 1.1\times2.40/6 = 6.34$): andesit. Lava 3 (66.50, 6.36; di antara 63
dan garis dari (69, 8) ke (77, 0), di bawah garis dari (63, 7) ke (69, 8)): dasit.

\textbf{5.} Data mentah: (50.4, 3.52), (58.6, 5.05), (65.3, 6.25): medan yang sama.
Perhitungan ulang memindahkan setiap titik paling banyak \qty{1.2}{\percent} silika dan tidak mengubah
apa pun di sini, berbeda dengan \cref{exo:b2:magmatic-series:2}.

\textbf{6.} \omterm{def:b2:magmatic-series:qapf}{Diagram QAPF}, dari modus yang terukur. Batuan plutonik yang berkomposisi
andesitik kaya akan plagioklas, dengan sedikit feldspar alkali atau kuarsa: batuan itu seharusnya
jatuh di medan 10 atau 10*, diorit atau diorit kuarsa (modusnya yang menentukan).

\noindent\textbf{Bagian II.}

\textbf{7.} Lava 1: garis pada 0.43 dan 1.47, $\ce{K2O} = 0.63$: K-sedang. Lava 2: 0.81
dan 2.28, 1.17: K-sedang. Lava 3: 1.13 dan 2.96, 1.88: K-sedang.

\textbf{8.} \ce{FeO_T} dan \ce{MgO} tanpa volatil: 8.67 dan 7.97; 6.18 dan 3.85; 4.07
dan 1.63. Lava 1: $3.55 + 8.67 + 7.97 = 20.19$; $A = 17.6$, $F = 42.9$, $M = 39.5$.
Lava 2: 15.15; 33.8, 40.8, 25.4. Lava 3: 12.06; 52.7, 33.7, 13.5.

\textbf{9.} Lava 1, $A = 17.6$: garis di antara (42.8, 17.1) dan (46.9, 18.1) bernilai
$F = 42.8 + 4.1\times0.5/1.0 = 44.9$; $42.9 < 44.9$, sisi kalk-alkali, dekat dengan
garis itu. Lava 2, $A = 33.8$: $F = 49.6 + 2.9\times1.2/5 = 50.3$; $40.8 < 50.3$,
kalk-alkali. Lava 3, $A = 52.7$: $F = 40.8 - 2.9\times2.7/5 = 39.2$; $33.7 < 39.2$,
kalk-alkali.

\textbf{10.} $\ce{FeO_T}/\ce{MgO}$: $8.67/7.97 = 1.09$, $6.18/3.85 = 1.61$, $4.07/1.63 =
2.50$. Rasio itu naik, seperti pada setiap basal yang terdiferensiasi, tetapi besi sendiri
turun dari 8.67 menjadi 4.07, dan $F$ pada segitiga turun dari 42.9 menjadi 33.7: tidak ada pengayaan
besi. Rangkaian toleiitik akan mula-mula naik ke arah F, seperti ferobasal pada
\cref{exo:b2:magmatic-series:5}.

\textbf{11.} \omterm{def:b2:magmatic-series:series}{Seri kalk-alkali}. Magnetit, yang mengkristal dari \omterm{def:b2:magma-emplacement:magma}{magma} yang teroksidasi,
mengeluarkan besi bersama magnesium; spesies teroksidasi itu berasal dari slab zona subduksi
yang melepaskan air.

\noindent\textbf{Bagian III.}

\textbf{12.} $\mathrm{Nb_N} = 2.0/0.658 = 3.04$, $\mathrm{La_N} = 7/0.648 = 10.80$,
$\mathrm{Th_N} = 1.1/0.0795 = 13.84$.

\textbf{13.} $3.04/10.80 = 0.281$, jauh di bawah median MORB (0.956) dan
\omterm{def:b2:elements-in-the-earth:morb-oib}{basal pulau samudra} (1.215), bahkan di bawah basal busur median (0.484): lembah
niobium yang dalam.

\textbf{14.} $13.84/3.04 = 4.55$; MORB median pada diagram laba-laba memiliki $2.64/6.08 =
0.43$. Torium, seperti unsur-unsur lain yang diperkaya dalam basal busur, ditambahkan oleh
slab yang menunjam.

\textbf{15.} \ce{TiO2} tanpa volatil $0.95\times1.00837 = 0.958$; $\mathrm{Ti} =
0.958\times5993 = \qty{5741}{ppm}$; Ti/100 = 57.4, Zr 70, $3\mathrm Y = 54$, jumlah 181.4:
31.6, 38.6, 29.8, di medan B, yang ditempati bersama oleh MORB, toleiit busur, dan basal kalk-alkali.
Medan itu menyingkirkan basal dalam-lempeng tetapi tidak memutuskan antara punggungan dan busur.
Diagram itu dibuat untuk basal: dalam andesit dan dasit, titanium telah disingkirkan
oleh oksida dan zirkonium dipekatkan oleh kristalisasi, dan medannya (lava 2
akan jatuh di C, lava 3 di luar semua medan) tidak akan berarti apa-apa.

\textbf{16.} \omterm{def:b2:magmatic-series:series}{Seri kalk-alkali} K-sedang, dengan lembah niobium yang dalam dan
$\mathrm{Th_N}/\mathrm{Nb_N}$ yang tinggi, pada gunung api yang meletuskan basal sampai dasit: gunung api itu
tumbuh di busur vulkanik di atas zona subduksi.

\noindent\textbf{Bagian IV.}

\textbf{17.} $C_L = C_0/F$ untuk unsur yang sepenuhnya inkompatibel: $F = 0.63/1.17 =
0.538$.

\textbf{18.} Rubidium $10/20 = 0.500$; torium $1.1/2.1 = 0.524$. (\omterm{def:b2:elements-in-the-earth:trace}{Unsur
jejak}, yang diberikan sebagaimana dianalisis, akan berubah kurang dari \qty{1}{\percent} oleh
perhitungan ulang.)

\textbf{19.} $C_L/C_0 = F^{D - 1}$: $125/70 = 1.786 = F^{-0.9}$, $F = 1.786^{-1/0.9} =
0.525$.

\textbf{20.} Dari kalium oksida, $F = 1.17/1.88 = 0.622$ dari lava 2 ke lava 3. Stronsium:
$380/480 = 0.792 = 0.622^{D - 1}$, sehingga $D - 1 = \ln 0.792/\ln 0.622 = -0.234/-0.474 =
0.49$ dan $D \approx 1.49$: stronsium telah menjadi kompatibel. Plagioklas, yang menyerap
stronsium, sedang mengkristal dalam jumlah besar.

\textbf{21.} Dari kalium oksida, $F = 0.63/1.88 = 0.335$ (rubidium $10/32 = 0.31$, torium
$1.1/3.4 = 0.32$): sekitar dua pertiga basal, $1 - 0.335 = 0.665$, telah mengkristal
sebelum dasit terbentuk.

\textbf{22.} Kalium oksida yang terasimilasi akan menaikkan \ce{K2O} dalam lava 2 melebihi yang
dihasilkan oleh kristalisasi saja, sehingga $F = 0.63/1.17$ akan terlalu rendah, dan massa
kristal terlalu tinggi.

\textbf{23.} $(0.538 + 0.500 + 0.524 + 0.525)/4 = 2.087/4 = 0.522$; kristal $1 - 0.522
= 0.478$ dari basal, yaitu \textbf{\qty[text-series-to-math=true]{48}{kg}
kristal per \qty{100}{kg} basal} untuk menghasilkan andesit itu.
\end{solution}
